\section{Boundaries and Non-Goals}

RCOUNT is most useful when its boundaries are explicit. It is not an election
management system, a voter registration database, a ballot marking device
verifier, an end-to-end voting protocol, or a replacement for risk-limiting
audits. It can store evidence relevant to those domains, but its core claim is
that public count packages can be replayed and checked.

The certification boundary is especially important. A canvassing board may
certify results after considering statutes, deadlines, cure rules, recount
requests, court orders, and evidence that cannot be reduced to arithmetic. RCOUNT
can show that a canvassed total equals a set of public summaries and that a
correction event exists. It cannot decide whether the board acted lawfully.

The privacy boundary is equally strict. A public artifact that lets a voter
prove their candidate choice can become a coercion receipt. RCOUNT therefore
treats inclusion proofs as privacy-gated sketches: they may show that an
anonymized token exists in accepted or counted evidence, but they must not carry
candidate selections or link voter identity, ballot style, and timestamp in a
way that reidentifies the ballot.

\begin{center}
\small
\begin{tabularx}{\textwidth}{lXX}
\toprule
Threat or claim & RCOUNT v0 position & Non-goal or warning \\
\midrule
Source tampering & Source bytes are hashed and checked against the source index. & Hashes do not prove the original source was truthful. \\
Parser substitution & Normalized records are package-hashed. & Parser diagnostics and independent parser disagreement are future adapter work. \\
Package truncation & Package content hash changes when normalized records are omitted. & Detached signatures are future work. \\
Equivocation & A transcript binds one package hash to one verification run. & Public log consistency is not yet implemented. \\
Small-cell disclosure & Proof sketches reject obvious linkable identity fields. & Rare write-ins, tiny precincts, ballot styles, and timestamps still require suppression policy. \\
Coercion receipts & Choice-bearing proof sketches fail the privacy gate. & RCOUNT must not let voters prove candidate choices. \\
Malware resistance & Out of scope. & RCOUNT is not a voting-system security proof. \\
\bottomrule
\end{tabularx}
\end{center}

The RPLAN boundary is optional. District vote aggregation is necessary when
precinct or reporting-unit totals must be projected onto districts, especially
when district assignments change across cycles. But count packages should not
depend on districting machinery unless the public claim being verified is a
district total. This keeps the base verifier small while allowing RCOUNT and
RPLAN to meet at a hashed assignment boundary.

\begin{center}
\small
\begin{tabularx}{\textwidth}{lX}
\toprule
RCOUNT can mechanically show & RCOUNT does not show \\
\midrule
Package records match their manifest hash. & Election officials acted lawfully. \\
Source bytes match the declared source index. & Voting-system software was uncompromised. \\
Contest and jurisdiction arithmetic reconciles. & Every eligible voter was registered or every ballot decision was correct. \\
Declared batches, lineage, and privacy gates pass. & Public proof sketches are a complete end-to-end voting protocol. \\
District totals replay from a specified RPLAN assignment. & The district plan is fair, lawful, or politically neutral. \\
\bottomrule
\end{tabularx}
\end{center}
